% !TeX program = pdflatex
% Build from the repository root; generated files stay in the ignored build tree:
% latexmk -pdf -outdir=build/docs/semismooth-contact docs/semismooth-rigid-contact-formulation.tex
% Technical documentation, not a submitted or published research article.
% The worked discretization and circle solver below retain the default P0 specialization.
% Configurable L2 ownership/assembly: docs/boundary-multiplier-space.md.
% 2026-09-09 boundary-space documentation update: no rebuild or new numerical tests.
\documentclass[preprint,11pt]{elsarticle}
\usepackage[T1]{fontenc}
\usepackage[utf8]{inputenc}
\usepackage{lmodern}
\usepackage{amsmath,amssymb,mathtools}
\usepackage{booktabs,array}
\usepackage{placeins}
\usepackage[margin=1in]{geometry}
\usepackage{microtype}
\usepackage{xurl}
\usepackage[hidelinks]{hyperref}
\hypersetup{pdftitle={Semismooth rigid contact for vector displacement: weak form and discretization}}
\journal{mfemMechanics technical documentation}
% Suppress elsarticle's submission-style first-page footer; retain page numbers.
\makeatletter
\let\ps@pfirstpage\ps@plain
\makeatother
\numberwithin{equation}{section}
\newcommand{\bu}{\boldsymbol{u}}
\newcommand{\bv}{\boldsymbol{v}}
\newcommand{\bw}{\boldsymbol{w}}
\newcommand{\bx}{\boldsymbol{x}}
\newcommand{\bX}{\boldsymbol{X}}
\newcommand{\bn}{\boldsymbol{n}}
\newcommand{\bN}{\boldsymbol{N}}
\newcommand{\bH}{\boldsymbol{H}}
\newcommand{\bI}{\boldsymbol{I}}
\newcommand{\bP}{\boldsymbol{P}}
\newcommand{\bF}{\boldsymbol{F}}
\newcommand{\bS}{\boldsymbol{S}}
\newcommand{\bC}{\mathbb{C}}
\newcommand{\Gc}{\Gamma_{c0}}
\newcommand{\ip}[2]{\left\langle #1,#2\right\rangle_{\Gc}}
\newcommand{\iq}[2]{\left\langle #1,#2\right\rangle_Q}
\newcommand{\pos}[1]{\left[#1\right]_+}
\newcommand{\Grad}{\operatorname{Grad}}
\newcommand{\eps}{\boldsymbol{\varepsilon}}
\newcommand{\Rbulk}{R_{\mathrm{bulk}}}
\newcommand{\Pibulk}{\Pi_{\mathrm{bulk}}}
\newcommand{\code}[1]{\texttt{\nolinkurl{#1}}}
\newcommand{\testname}[1]{\par\smallskip\noindent\code{#1}\par\nobreak}
\setlength{\emergencystretch}{2em}

\begin{document}
\begin{frontmatter}
\title{Semismooth rigid contact for vector displacement:\\
gap kinematics, mixed weak form, and finite element discretization\\[3pt]
\large A formulation note for mfemMechanics technical documentation}
% Intentionally no research-author or affiliation attribution.
\begin{abstract}
This technical documentation derives the displacement--multiplier equations
implemented by \texttt{SemismoothRigidContactOperator} in mfemMechanics.
It is an implementation-oriented adaptation of the contact algebra in
Burman, Hansbo, and Larson's augmented Lagrangian formulation, not a claim of
publication, submission, or a new stability theorem. The starting point is
the vector-displacement counterpart of Eq.~(2.3) in the authors'
12 September 2016 arXiv v1 PDF. A positive signed gap and a positive
compressive nominal multiplier are used throughout. The initial gap is
separated from displacement-induced closure, exactly for a plane and through
an integral remainder for a curved rigid wall. A reference-boundary saddle
functional then gives the mixed residual, default boundary piecewise-constant
multiplier stabilization, and all four generalized Jacobian blocks,
including the signed-distance curvature term. The finite element and
quadrature equations are developed explicitly, followed by a hand-integrated
two-node edge, essential-constraint handling, primal-field extraction, and
the monolithic Newton solve. Named source tests provide a verification map;
no new numerical experiments are reported. In particular, discrete face
moments and algebraic convergence are not presented as pointwise
nonpenetration or as evidence of mesh-independent stability.
\end{abstract}
\begin{keyword}
Rigid contact \sep Vector elasticity \sep Signed distance \sep
Augmented Lagrangian \sep Semismooth Newton method \sep Boundary P0 multiplier
\end{keyword}
\end{frontmatter}

\section{The vector counterpart of the source mixed equations}
\label{sec:lead}

Let $\bu$ be the body's \emph{vector displacement}, $\bv$ a displacement
test field, $g(\bu;\bX)$ the signed distance from the displaced boundary
point to a fixed rigid wall, and $\lambda$ a scalar multiplier positive
in compression. Define the augmented nominal pressure
\begin{equation}
 p(\bu,\lambda;\bX)=\pos{\lambda-\frac{g(\bu;\bX)}{\gamma}},
 \qquad \gamma>0,
 \qquad \bn(\bu;\bX)=\nabla_x d_w(\bX+\bu(\bX)).
 \label{eq:pressure}
\end{equation}
Here $d_w$ is the wall's signed-distance function, positive in the admissible
region, and $\gamma$ is a compliance with dimensions length/stress.
With the reference-boundary product
$\ip{a}{b}=\int_{\Gc}ab\,dA_0$, the desired counterpart of the
source PDF's Eq.~(2.3) is the following \emph{unstabilized} pair:
\begin{equation}
 \boxed{\begin{aligned}
  \Rbulk(\bu;\bv)-\ip{p}{\bn\cdot\bv}&=0,
       &&\text{for every admissible }\bv,\\
  \ip{\lambda-p}{\mu}&=0,
       &&\text{for every multiplier test }\mu.
 \end{aligned}}
 \label{eq:counterpart}
\end{equation}
The bulk residual is internal virtual work minus prescribed external
virtual work. Thus the actual contact force on the body has positive
direction $p\bn$, whereas its contribution to the residual is negative.
The second row of \eqref{eq:counterpart} has the normalization of the
source Eq.~(2.3). The code instead multiplies this row by $-\gamma$,
giving $\gamma\ip{p-\lambda}{\mu}=0$ before stabilization.
This normalization is essential when comparing tangents or adding the
negative multiplier jump term.

Equation~\eqref{eq:counterpart} is a smooth-formal mechanical statement,
not an assertion that a continuous $L^2$ multiplier formulation is
well posed. The precise problem documented here is the finite-dimensional,
quadrature-defined system in Section~\ref{sec:discrete}, with
$\lambda_h\in M_h\subset L^2(\Gc)$. The preceding boundary integrals
presuppose enough regularity to perform the indicated calculations.
In particular, a scalar contact expression is being transferred to a
vector bulk model; the scalar and vector field equations are not identified.

\paragraph{Provenance and reading scope.}
The source is Burman, Hansbo, and Larson,
\emph{Augmented Lagrangian finite element methods for contact problems}
\cite{BHLv1}. Source verification for this documentation
identifies Zotero record \code{ZAV78EWP} and the PDF's first-page mark
``arXiv:1609.03326v1'', dated 12 September 2016, in a 26-page document.
Reading PDF pages 1--4 and visually checking pages 2 and 4 verifies Section~2,
Eqs.~(2.1)--(2.3), (2.7), and (2.10), used below.
These are \emph{v1 PDF locators}, not journal equation numbers.
The separately rendered HTML numbers the corresponding expressions
(3)--(5), (9), and (12). The later journal metadata is recorded in
\cite{BHLjournal}; no journal-numbering equivalence is presumed.
The inspected source is the working tree based on \code{bf40b5f}, including
its pre-existing uncommitted changes. Current contact source, the circle
example, and the named tests were inspected. Gap decompositions, vector chain rules, matrix
blocks, and the edge calculation below are derivations adapted to the
repository, not theorems quoted from the scalar paper. The earlier \LaTeX{}
manuscript was compiled and its equation layout inspected; the boundary-space
documentation update has not been rebuilt, and reports no new C++ tests or
numerical experiments.

\paragraph{Current boundary-space API.}
The worked discretization below retains the default P0 specialization.
The current operator instead borrows a \code{BoundaryMultiplierSpace}, whose
null collection selects P0 and whose configurable collection accepts scalar,
value-mapped discontinuous L2 elements. The argument after the obstacle is
that boundary object, not the former attribute array; public contact accessors
are unchanged. Ownership, basis-aware assembly, P1/Q1 regression constructions,
and limits are documented in \code{docs/boundary-multiplier-space.md}.
This extension is not a continuous/IGA implementation or new stability evidence.

\section{Reference configuration, fields, and mechanical scope}
\label{sec:scope}

\subsection{Domain and admissible variations}
Let $\Omega_0\subset\mathbb{R}^d$, $d\in\{2,3\}$, be the fixed
reference body, and let $\Gc\subset\partial\Omega_0$ be the selected
exterior candidate contact boundary. Candidate does not mean active:
only part of this boundary may carry pressure in a given iterate.
The reference mesh and this selection remain fixed during a solve.
At a material point $\bX$, the placement is
\begin{equation}
 \bx(\bX)=\bX+\bu(\bX),\qquad
 \bF=\Grad\bx=\bI+\Grad\bu.
 \label{eq:placement}
\end{equation}
The mechanics considered here is quasi-static. The operator evaluates rigid
contact only on $\Gc$, not in the body
volume. It neither searches between two deformable meshes nor constructs
slave/master pairs for mesh-to-mesh contact.

Component-wise prescribed displacements are represented by sets
$\Gamma_D^{(i)}\subset\partial\Omega_0$ and data $u_D^{(i)}$:
\begin{align}
 V_D&=\{\bu\in[H^1(\Omega_0)]^d:
       u_i=u_D^{(i)}\text{ on }\Gamma_D^{(i)},\ i=1,\ldots,d\},\\
 V_0&=\{\bv\in[H^1(\Omega_0)]^d:
       v_i=0\text{ on }\Gamma_D^{(i)},\ i=1,\ldots,d\}.
 \label{eq:spaces}
\end{align}
$V_D$ is affine, while $V_0$ is the associated homogeneous variation
space. Sufficiently regular members are understood when signed-distance
derivatives are evaluated pointwise. The total displacement belongs to
$V_D$; a Newton correction $\Delta\bu$ and a test $\bv$ belong to $V_0$.
They play different roles even though their finite element bases coincide.
Nonzero prescribed displacement must remain in \eqref{eq:placement} when
evaluating the gap.

\subsection{Signed distance, normals, and dimensions}
Assume $d_w$ is $C^2$ in a neighborhood of all evaluated trial positions.
Define
\begin{equation}
 g(\bu;\bX)=d_w(\bX+\bu(\bX)),\qquad
 \bn=\nabla_x d_w,\qquad \bH=\nabla_x^2 d_w,
 \qquad \|\bn\|=1.
 \label{eq:distance}
\end{equation}
The wall-side normal $\bn$ points into the admissible region. It is not
the body outward reference normal $\bN$. Even the current body outward
normal equals $-\bn$ only where the two surfaces are appropriately
aligned. Confusing these normals changes the contact force sign.

The ideal local unilateral conditions are
\begin{equation}
 g\geq0,\qquad \lambda\geq0,\qquad \lambda g=0.
 \label{eq:complementarity}
\end{equation}
Frictionlessness means zero traction in every direction orthogonal to the
\emph{local} wall normal: $(\bI-\bn\bn^T)(p\bn)=0$. On a curved wall,
this traction can have a nonzero component along a fixed global horizontal
direction, or along a reference-body tangent. Such a component is not a
friction law. There is no tangential penalty, adhesion, slip history, or
contact rate law in this formulation.

\begin{table}[htbp]
\centering
\caption{Dimensions and roles. Stress denotes force per area; the 2D model
uses the corresponding quantities per unit out-of-plane thickness.}
\label{tab:units}
\begin{tabular}{@{}p{0.26\linewidth}p{0.43\linewidth}p{0.22\linewidth}@{}}
\toprule
Quantity & Role & Dimension\\
\midrule
$\bX,\bx,\bu,\bv,g$ & Position, displacement, variation, gap & length\\
$\lambda,p,\mu$ & Nominal multiplier, pressure, variation & stress\\
$\bn,\bN$ & Unit normals with distinct meanings & dimensionless\\
$\bH$ & Spatial signed-distance Hessian & inverse length\\
$\gamma$ & Fixed augmentation compliance & length/stress\\
$\delta$ & Fixed jump factor, $\delta\geq0$ & dimensionless\\
$h_F$ & Reference contact mesh scale & length\\
$dA_0,w_q$ & Reference boundary measure, quadrature weight & length$^{d-1}$\\
\bottomrule
\end{tabular}
\end{table}

\subsection{Bulk mechanics and nominal traction}
For a hyperelastic body with stored-energy density $W(\bF)$ per reference
volume and dead external loads, take
\begin{align}
 \Pibulk(\bu)&=\int_{\Omega_0}W(\bF)\,dV_0-\ell(\bu),\\
 \Rbulk(\bu;\bv)&=D\Pibulk(\bu)[\bv]
 =\int_{\Omega_0}\bP:\Grad\bv\,dV_0-\ell(\bv),
 \qquad \bP=\frac{\partial W}{\partial\bF}.
 \label{eq:bulk}
\end{align}
For example, $\ell(\bv)$ includes reference body-force and prescribed
reference-traction integrals. These are dead loads; follower loads would
need their own consistent derivatives. The stress in
\eqref{eq:bulk} is first Piola--Kirchhoff stress. If the constitutive API
returns second Piola--Kirchhoff stress $\bS$, then $\bP=\bF\bS$;
the Cauchy stress is $\boldsymbol{\sigma}=J^{-1}\bF\bS\bF^T$, with
$J=\det\bF>0$. The primal formulation, not the contact operator, must
enforce admissible volume kinematics and supply its material and geometric
tangent terms.

Small-strain elasticity specializes this notation to
\begin{equation}
 \eps(\bu)=\tfrac12(\Grad\bu+\Grad\bu^T),\qquad
 a(\bu,\bv)=\int_{\Omega_0}\eps(\bv):\bC:\eps(\bu)\,dV_0,
 \qquad \Rbulk=a(\bu,\bv)-\ell(\bv).
 \label{eq:smallstrain}
\end{equation}
The repository's ordinary 2D solid convention is plane strain, not an
implicit plane-stress reduction. The circle example uses small-strain
bulk elasticity but the exact nonlinear distance in \eqref{eq:distance}.
This is a hybrid modeling choice: small bulk strains do not make a curved
gap affine, nor do exact gap kinematics make the bulk law finite-strain
objective for arbitrary large rotations.

The applied contact traction $p\bn$ is per reference boundary measure.
If the deformed surface is regular and $j_s=dA/dA_0$ is its surface ratio,
the same force written per current area has magnitude
\begin{equation}
 p_{\mathrm{current}}=\frac{p}{j_s},\qquad
 j_s=J\|\bF^{-T}\bN\|.
 \label{eq:surfacepressure}
\end{equation}
The second expression assumes the finite-deformation surface map is
applicable. Reported $p$ is therefore not automatically a physical
current-area pressure. The assembly keeps $dA_0$ fixed and takes no
displacement derivative of it; a current-area contact potential would be
a different formulation with additional derivatives.

\section{Initial gap and displacement-induced closure}
\label{sec:gap}

\subsection{Exact decomposition for a plane}
Define the initial clearance and initial wall normal by
\begin{equation}
 g_0(\bX)=d_w(\bX),\qquad \bn_0(\bX)=\nabla_x d_w(\bX).
 \label{eq:initialgap}
\end{equation}
For a plane through $\bx_w$ with unit admissible normal $\bn_0$,
$d_w(\bx)=(\bx-\bx_w)\cdot\bn_0$. Consequently,
\begin{equation}
 \boxed{g(\bu;\bX)=g_0(\bX)+\bn_0\cdot\bu(\bX)},\qquad
 \eta(\bu;\bX):=-g=-g_0-\bn_0\cdot\bu.
 \label{eq:planegap}
\end{equation}
This identity is exact, not a first-order approximation. The signed
penetration variable $\eta$ is positive inside the obstacle; its positive
part is the penetration depth. The displacement trace is still a vector.
Even for a plane, $\eta$ is not just a normal displacement when $g_0\ne0$.

For a horizontal floor $y=y_0$ with the body above it, the admissible
normal is $\bn_0=\boldsymbol{e}_y$, so
\begin{equation}
 g_0=X_y-y_0,\qquad g=X_y-y_0+u_y,\qquad
 \eta=y_0-X_y-u_y.
 \label{eq:floor}
\end{equation}
An initially separated point first touches when $u_y=-g_0$. Further
downward motion makes $g<0$, and the contact force points upward.
The body outward normal on an aligned lower horizontal boundary points
downward, which explains the minus sign relating that normal to $\bn_0$.

For small-strain bulk elasticity and this planar wall, substituting
\eqref{eq:planegap} into \eqref{eq:counterpart} makes the vector counterpart
of source Eq.~(2.3) completely explicit:
\begin{equation}
 \boxed{\begin{aligned}
 a(\bu,\bv)
 -\left\langle
    \pos{\lambda-\frac{g_0+\bn_0\cdot\bu}{\gamma}},
    \bn_0\cdot\bv
  \right\rangle_{\Gc}&=\ell(\bv),\\
 \left\langle
    \lambda-\pos{\lambda-\frac{g_0+\bn_0\cdot\bu}{\gamma}},
    \mu
  \right\rangle_{\Gc}&=0.
 \end{aligned}}
 \label{eq:expandedplane}
\end{equation}
Here $\bu\in V_D$, $\bv\in V_0$, and the second row is still the
\emph{unstabilized, normalized} multiplier equation. In particular,
$g_0$ is prescribed geometric data, so its variation is zero; the wall-normal
displacement $\bn_0\cdot\bu$ is the part varied in Newton. The implemented
multiplier row is obtained by multiplying this second equation by
$-\gamma$ and then adding $-s_h(\lambda_h,\mu_h)$ in the discrete problem.

\subsection{Exact curved-wall identities and their restrictions}
Fix $\bX$ and $\bu(\bX)$, and consider the scalar function
$f(t)=d_w(\bX+t\bu)$ for $0\leq t\leq1$.
If this entire segment lies in a smooth signed-distance neighborhood, the
fundamental theorem of calculus gives
\begin{equation}
 g(\bu;\bX)=g_0(\bX)
   +\int_0^1\bn(\bX+t\bu)\cdot\bu\,dt.
 \label{eq:gapintegral}
\end{equation}
Here $\bn(\bx)$ denotes the spatial gradient of $d_w$ evaluated at the
displayed spatial argument. A second integration gives the exact Taylor
formula with integral remainder,
\begin{equation}
 \boxed{g(\bu;\bX)=g_0+\bn_0\cdot\bu
 +\int_0^1(1-t)\,\bu^T\bH(\bX+t\bu)\bu\,dt.}
 \label{eq:gapremainder}
\end{equation}
Indeed, $f'(t)=\bn(\bX+t\bu)\cdot\bu$ and
$f''(t)=\bu^T\bH(\bX+t\bu)\bu$. These intermediate derivatives
show exactly where curvature enters the closure. If
$\|\bH\|\leq C_H$ on the segment, the remainder is bounded in magnitude
by $C_H\|\bu\|^2/2$. Thus $g_0+\bn_0\cdot\bu$ is a first-order
approximation under that local bound, not the nonlinear law used by the
operator.

In particular, the expression $g_0+\bn(\bX+\bu)\cdot\bu$ is
\emph{not} an exact curved-wall decomposition. The end-point normal is
neither the path average in \eqref{eq:gapintegral} nor the initial
normal plus the correct integral remainder. Nor should the path assumption
be imposed unnecessarily on assembly: the operator evaluates the distance
at the current quadrature position. Its local Newton derivative requires
smoothness near that position, whereas \eqref{eq:gapintegral} additionally
requires smoothness along the whole reference-to-current segment.

\subsection{Linearization at the current iterate}
At a current displacement $\bu$, a small increment $\Delta\bu$ gives
\begin{align}
 g(\bu+\Delta\bu;\bX)
  &=g(\bu;\bX)+\bn(\bu;\bX)\cdot\Delta\bu
       +O(\|\Delta\bu\|^2),\label{eq:gapnewton}\\
 Dg(\bu)[\bv]&=\bn\cdot\bv,\qquad
 D\bn(\bu)[\bw]=\bH\bw,\qquad
 D^2g(\bu)[\bv,\bw]=\bv^T\bH\bw.
 \label{eq:gapderivatives}
\end{align}
These are current-iterate derivatives of the exact distance law, not an
instruction to accumulate a linearized gap from accepted load increments.
The implementation recomputes $g$ from the current total placement.

For a circle in 2D or sphere in 3D, with center $\boldsymbol{c}$,
radius $R>0$, and admissible exterior, let
$\boldsymbol{r}=\bX+\bu-\boldsymbol{c}$ and $r=\|\boldsymbol{r}\|$.
Then
\begin{equation}
 g=r-R,\qquad \bn=\frac{\boldsymbol{r}}r,\qquad
 \bH=\frac{\bI-\bn\bn^T}{r},\qquad r>0.
 \label{eq:sphere}
\end{equation}
The center is excluded because the distance is not differentiable there.
General walls likewise require exclusion of cut loci and nonsmooth
closest-point regions. The identities $\bH^T=\bH$ and $\bH\bn=0$
follow for a smooth signed distance and are checked by the code, together
with finiteness and unit gradient. These necessary identities do not
independently prove that a custom obstacle's returned derivatives
differentiate its returned gap.

\section{From scalar contact algebra to the vector saddle functional}
\label{sec:adaptation}

\subsection{Sign and test-function map}
Use a subscript $B$ for the scalar source variables. The verified source
PDF Eq.~(2.1) and mixed pair (2.3) read, with its contact product understood,
\begin{align}
 \lambda_B&=-\frac1\gamma\pos{u_B-\gamma\lambda_B},
 \label{eq:sourceprojection}\\
 a_B(u_B,v_B)+\left\langle\frac1\gamma
       \pos{u_B-\gamma\lambda_B},v_B\right\rangle_C
       &=(f,v_B),\label{eq:sourcerow1}\\
 \left\langle\lambda_B+\frac1\gamma
       \pos{u_B-\gamma\lambda_B},\mu_B\right\rangle_C&=0.
 \label{eq:sourcerow2}
\end{align}
These expressions are transcriptions with renamed symbols
\cite[Section 2, PDF Eqs.~(2.1), (2.3)]{BHLv1}.
Only their \emph{contact terms} are adapted by
\begin{equation}
 u_B=-g=\eta,\qquad \lambda_B=-\lambda,\qquad
 v_B=D(-g)[\bv]=-\bn\cdot\bv,\qquad \mu_B=-\mu.
 \label{eq:sourcemap}
\end{equation}
The test map is a chain-rule variation of the gap map. For curvature it
depends on the current displacement, so differentiating it again cannot
be omitted. The source bulk bilinear form $a_B$ is replaced by the
mechanical residual of Section~\ref{sec:scope}; it is not transformed into
elasticity by \eqref{eq:sourcemap}.

Positive homogeneity for $\gamma>0$ gives
\begin{equation}
 P_B^+:=\pos{u_B-\gamma\lambda_B}
       =\pos{-g+\gamma\lambda}=\gamma p.
 \label{eq:positivehomogeneity}
\end{equation}
Substitution into the two source contact rows produces
$-\ip{p}{\bn\cdot\bv}$ and $\ip{\lambda-p}{\mu}$,
respectively, establishing \eqref{eq:counterpart} with the adopted
bulk residual.

Pointwise, $\lambda=p$ is equivalent to \eqref{eq:complementarity}:
if $\lambda>0$, the positive-part branch gives $g=0$; if $\lambda=0$,
the map requires $g\geq0$. Conversely, either complementary case satisfies
the map. This is a pointwise algebraic equivalence, not a license to
replace $p$ by $\lambda_h$ inside a discrete stabilized displacement row.
Newton iterates of $\lambda_h$ are not clipped to be nonnegative; the
projection is applied to $\lambda_h-g/\gamma$ to compute $p$.

\subsection{First variation and multiplier-row normalization}
Adapting the contact part of source PDF Eq.~(2.2), define formally
\begin{equation}
 \mathcal{L}_\gamma(\bu,\lambda)=\Pibulk(\bu)
       +\frac\gamma2\bigl(\ip{p}{p}-\ip{\lambda}{\lambda}\bigr).
 \label{eq:functional}
\end{equation}
The source term is $\|P_B^+\|_C^2/(2\gamma)
-\gamma\|\lambda_B\|_C^2/2$; \eqref{eq:functional} follows from
\eqref{eq:sourcemap}--\eqref{eq:positivehomogeneity}
\cite[Section 2, PDF Eq.~(2.2)]{BHLv1}.
The scalar function $z\mapsto\tfrac12\pos{z}^2$ is continuously
differentiable at zero, with derivative $\pos{z}$. Consequently its first
variation needs no arbitrary derivative choice at the kink:
\begin{align}
 D\mathcal{L}_\gamma[\bv,\mu]
 &=\Rbulk(\bu;\bv)
   +\gamma\ip{p}{\mu-\gamma^{-1}\bn\cdot\bv}
   -\gamma\ip{\lambda}{\mu}\notag\\
 &=\Rbulk(\bu;\bv)-\ip{p}{\bn\cdot\bv}
                 +\gamma\ip{p-\lambda}{\mu}.
 \label{eq:firstvariation}
\end{align}
The multiplier first variation is $-\gamma$ times the normalized second
row of \eqref{eq:counterpart}. Although these unstabilized equations have
the same zeros, only \eqref{eq:firstvariation} has the saddle-functional
normalization used in the code.

The stabilized contact term of source Formulation~1 is
\begin{equation}
 b_B=\left\langle\gamma^{-1}P_B^+,v_B-\gamma\mu_B\right\rangle_C
       -\left\langle\gamma\mu_B,\lambda_B\right\rangle_C
       -s_B(\lambda_B,\mu_B),
 \label{eq:sourceb}
\end{equation}
as verified in PDF Eq.~(2.10) \cite{BHLv1}. The simultaneous sign changes
in multiplier trial and test leave the bilinear jump product unchanged.
Thus \eqref{eq:sourceb} gives exactly the contact part of
\eqref{eq:firstvariation} minus $s_h(\lambda_h,\mu_h)$.
Adding $-s_h(\lambda_h,\lambda_h)/2$ to the discrete functional yields
this term directly. There is no extra, ad hoc $\gamma$ rescaling of the
assembled stabilization matrix. A convention with a normalized second
row would have to rescale \emph{that entire row}, including its jump term.

The mixed formulation retains projected $p$ in the primal row. It is not
the alternative formulation obtained by putting the multiplier itself
there. Moreover, \eqref{eq:functional} is a saddle functional rather than
a positive contact energy to minimize jointly over both unknowns.
The implementation assembles derivatives but exposes no mixed energy
evaluation API or proved energy-based globalization.

\section{Finite element problem and reference quadrature}
\label{sec:discrete}

\subsection{Spaces, traces, and coefficient ordering}
Let $\mathcal{T}_h$ be the conforming reference body mesh and
$\mathcal{C}_h$ its selected boundary mesh. The H1 displacement spaces
$V_{D,h}\subset V_D$ and $V_{0,h}\subset V_0$ use the same polynomial
degree and the same essential constraints. The default P0 multiplier
space developed here is
\begin{equation}
 M_h=\{\lambda_h\in L^2(\Gc):
             \lambda_h|_e\in\mathbb{P}_0(e),\ e\in\mathcal{C}_h\}.
 \label{eq:p0space}
\end{equation}
There is exactly one independent multiplier coefficient per selected
boundary face. Writing $M_e$ for its indicator basis,
\begin{equation}
 \bu_h=\sum_a N_a\boldsymbol{u}_a,\qquad
 \bv_h=\sum_a N_a\boldsymbol{v}_a,\qquad
 \lambda_h=\sum_e M_e\Lambda_e,\qquad
 \mu_h=\sum_f M_f\mathcal{M}_f.
 \label{eq:febases}
\end{equation}
The vector test coefficients vanish on essential components; total
displacement coefficients retain prescribed values. $\Delta\boldsymbol{u}_a$
are Newton corrections, not test coefficients. Neither $\bv_h$ nor
$\mu_h$ is an additional unknown field.

On a particular contact face with $n_b$ scalar H1 trace functions, the
component-major local vector and its evaluation matrix are
\begin{equation}
 U_e=\begin{bmatrix}
 u_{1,1}&\cdots&u_{n_b,1}&u_{1,2}&\cdots&u_{n_b,d}
 \end{bmatrix}^{T},\qquad
 \mathsf N_q=\bI_d\otimes
      \begin{bmatrix}N_1(X_q)&\cdots&N_{n_b}(X_q)\end{bmatrix}.
 \label{eq:evalmatrix}
\end{equation}
Here $\mathsf N_q\in\mathbb{R}^{d\times dn_b}$ and
$\bu_h(X_q)=\mathsf N_qU_e$. The scalar multiplier evaluation matrix
is $\mathsf M_q\in\mathbb{R}^{1\times n_m}$ with $n_m=1$ and
$\mathsf M_q=[1]$ locally for P0. The current configurable L2 implementation
uses the same formulas with the full local basis row and its actual $n_m$;
the face-moment and edge examples below remain P0-specific.

The code obtains boundary vector DOFs and boundary H1 shapes, not all
adjacent volume shape functions for assembly. Volume DOFs with zero trace
therefore have no contact contribution. MFEM's local component-major
layout is respected independently of the configured global vector ordering;
global assembly uses the space's DOF maps. Contact uses trace values only:
no surface derivative of a multiplier basis and no body displacement
gradient occurs in its local residual. Gradients needed by the bulk law
remain the responsibility of volume assembly.

\subsection{Fixed reference quadrature and nonlinear evaluation}
For reference quadrature points $\widehat{\xi}_q$ and weights
$\widehat w_q$ on a selected face, define
\begin{align}
 w_q&=\widehat w_q J_{\Gamma_0}(\widehat{\xi}_q),\qquad
 \iq{a}{b}=\sum_{e\in\mathcal C_h}\sum_{q\in e}w_q a_qb_q,
 \label{eq:quadrature}\\
 \bx_q&=\bX_q+\mathsf N_q U_e,\qquad
 g_q=d_w(\bx_q),\qquad \lambda_q=\mathsf M_q\Lambda_e,\notag\\
 z_q&=\lambda_q-g_q/\gamma,\qquad
 p_q=\max(0,z_q),\qquad
 \chi_q=\begin{cases}1,&z_q\geq0,\\0,&z_q<0.\end{cases}
 \label{eq:pointkernel}
\end{align}
$J_{\Gamma_0}$ is the reference surface metric. In code $w_q$ is
\code{integrationPoint.weight * boundaryTransformation->Weight()}.
The integration point is set before geometry evaluation, and inverse
boundary/mesh-face orientation aligns trace shapes with the selected
mesh-face-coordinate rule, including nonsymmetric custom rules.
The metric enters exactly once and has no derivative with respect to $U_e$.

Displacement is interpolated first; distance, normal, Hessian, positive
part, and activity are then evaluated at \emph{each} quadrature point.
The algorithm does not interpolate a nodally computed contact pressure.
Even though $\lambda_h$ is constant on a face, $g_q$, $p_q$, and $\chi_q$
need not be constant there. Interpolation, averaging, and the positive-part
operation cannot in general be interchanged.

Residual and Jacobian must use the same fixed rule. The default rule order
is based on adjacent-volume geometry and the maximum of displacement and
multiplier orders (retaining the P0 policy), with the
code's Pk increment; it is not a guarantee of accurate integration through
a curved active-set boundary. A custom rule is borrowed and must remain
alive. Changing samples during Newton changes the discrete residual being
solved, rather than merely improving its linearization.

\subsection{Boundary jump stabilization}
Let $\mathcal F_h^c$ contain interior interfaces of the contact submesh,
each counted once. For neighboring contact faces $e,f$ meeting at $F$,
choose consistently oriented jumps
$[\lambda_h]_F=\Lambda_e-\Lambda_f$ and
$[\mu_h]_F=\mathcal M_e-\mathcal M_f$. Adapting the source PDF
Eq.~(2.7), p.~4 \cite{BHLv1}, the code uses
\begin{align}
 s_h(\lambda_h,\mu_h)
 &=\sum_{F\in\mathcal F_h^c}\delta\gamma h_F
        \int_F[\lambda_h]_F[\mu_h]_F\,ds_0,
 \label{eq:jump}\\
 h_F&=\frac{|e|^{1/d_c}+|f|^{1/d_c}}2,\qquad d_c=d-1.
 \label{eq:hface}
\end{align}
The local area-root length scale in \eqref{eq:hface} is the repository's
specific choice. The measures are reference quantities. When $d=2$, an
interface is a contact-mesh vertex with measure $|F|=1$; when $d=3$, it
is an edge with integrated reference length. It is not a volume face or
the area of a whole contact element.
In the implementation, geometric measures are numerically evaluated and
cached; in particular, a curved interface length is evaluated by quadrature,
not generally exactly. The P0 reduction below is exact algebra for those
computed coefficients, not a claim of exact curved-geometry integration.

For P0, the interface contribution reduces exactly to
\begin{equation}
 c_F(\Lambda_e-\Lambda_f)(\mathcal M_e-\mathcal M_f),\qquad
 c_F=\delta\gamma h_F|F|,\qquad
 S_F=c_F\begin{bmatrix}1&-1\\-1&1\end{bmatrix}.
 \label{eq:jumpmatrix}
\end{equation}
Globally, $s_h(\lambda_h,\mu_h)=\mathcal M^T S\Lambda$.
$S$ is positive semidefinite for nonnegative $c_F$, has zero row sums,
and leaves constant multipliers on each connected contact component
unchanged. Reversing both jumps changes nothing. The residual contains
$-S\Lambda$ and its multiplier derivative contains $-S$: the negative
sign belongs to the saddle block, not to a negative body-stiffness penalty.
The term acts across all selected interior interfaces, including
active/inactive neighbors. $\delta=0$ disables it; the default $\delta=1$
is an implementation choice, not a general stability guarantee.

\subsection{Final mixed discrete weak form}
The quadrature-defined saddle functional is
\begin{equation}
 \mathcal L_{\gamma,h}(\bu_h,\lambda_h)=\Pi_{\mathrm{bulk},h}(\bu_h)
 +\frac\gamma2\bigl(\iq{p_h}{p_h}-\iq{\lambda_h}{\lambda_h}\bigr)
 -\frac12s_h(\lambda_h,\lambda_h).
 \label{eq:discretefunctional}
\end{equation}
For a conservative primal model its stationarity is the implemented
discrete problem: find $(\bu_h,\lambda_h)\in V_{D,h}\times M_h$ such that
\begin{equation}
 \boxed{\begin{aligned}
 r_u(\bv_h)&=R_{\mathrm{bulk},h}(\bu_h;\bv_h)
                   -\iq{p_h}{\bn_h\cdot\bv_h}=0,\\
 r_\lambda(\mu_h)&=\gamma\iq{p_h-\lambda_h}{\mu_h}
                   -s_h(\lambda_h,\mu_h)=0,
 \end{aligned}\quad
 \begin{gathered}
   \forall\,\bv_h\in V_{0,h},\\
   \forall\,\mu_h\in M_h.
 \end{gathered}}
 \label{eq:finalweak}
\end{equation}
The bulk residual may also be supplied by a more general primal operator;
then the contact functional still explains the contact rows, without
asserting a potential for the entire primal problem. $\gamma$ and
$\delta$ remain fixed parameters in every derivative.

Testing by unit coefficient directions gives the unscaled assembled
vectors, with face-to-global insertion understood:
\begin{align}
 r_U&=r_{\mathrm{bulk}}-\sum_{e,q}w_qp_q\mathsf N_q^T\bn_q,
 \label{eq:vectorru}\\
 r_\Lambda&=\gamma\sum_{e,q}w_q\mathsf M_q^T
                          (p_q-\mathsf M_q\Lambda_e)-S\Lambda.
 \label{eq:vectorrl}
\end{align}
The contact portions have local dimensions $dn_b$ and $n_m=1$,
respectively. Displacement residual coefficients have units
stress times length$^{d-1}$, whereas multiplier residual coefficients
have units length$^d$. Pairing them with displacement and multiplier
variations produces energy in 3D, or energy per thickness in 2D, in both
cases. This distinction motivates explicit row scaling rather than
comparing unlike raw coefficient norms.

\section{Generalized linearization and the four matrix blocks}
\label{sec:tangent}

\subsection{Directional derivation}
For an increment $(\Delta\bu_h,\Delta\lambda_h)$ at a fixed smooth
geometry evaluation, the selected generalized pressure derivative is
\begin{equation}
 \Delta g=\bn\cdot\Delta\bu_h,\qquad
 \Delta\bn=\bH\Delta\bu_h,\qquad
 \Delta p=\chi\left(\Delta\lambda_h-
                       \gamma^{-1}\bn\cdot\Delta\bu_h\right).
 \label{eq:pointderivative}
\end{equation}
Away from $z=0$, this is the ordinary derivative on the current branch.
At $z=0$, the positive part has generalized derivative interval $[0,1]$;
the code selects $\chi=1$, the active-side value. It is not the ordinary
directional derivative for every direction crossing the switching point.
The squared positive part in the potential is differentiable, but its
generalized Hessian still requires this choice.

Keep test fields fixed while differentiating the residual. The product
rule applied to $-p\bn\cdot\bv_h$ gives two terms, not one:
\begin{align}
 Dr_u[\Delta\bu_h,\Delta\lambda_h;\bv_h]
 &=D R_{\mathrm{bulk},h}[\Delta\bu_h;\bv_h]
   +\iq{\gamma^{-1}\chi\,\bn\cdot\Delta\bu_h}
        {\bn\cdot\bv_h}\notag\\
 &\quad-\iq{\chi\Delta\lambda_h}{\bn\cdot\bv_h}
   -\iq{p}{\bv_h^T\bH\Delta\bu_h},
 \label{eq:directionru}\\
 Dr_\lambda[\Delta\bu_h,\Delta\lambda_h;\mu_h]
 &=-\iq{\chi\,\bn\cdot\Delta\bu_h}{\mu_h}
    +\gamma\iq{(\chi-1)\Delta\lambda_h}{\mu_h}
    -s_h(\Delta\lambda_h,\mu_h).
 \label{eq:directionrl}
\end{align}
The term involving $\chi/\gamma$ comes from pressure augmentation; the
term involving $-p\bH$ comes from rotation of the wall normal. In
particular, the displacement block is not merely the bulk tangent of
an ordinary Lagrange-multiplier method.

\subsection{Coefficient matrices}
Substitution of the trace evaluation matrices into
\eqref{eq:directionru}--\eqref{eq:directionrl} yields
\begin{align}
 K_{11}&=K_{\mathrm{bulk}}+K_{uu}^{c},\qquad
 K_{uu}^{c}=\sum_{e,q}w_q\mathsf N_q^T
       \left(\frac{\chi_q}{\gamma}\bn_q\bn_q^T-p_q\bH_q\right)
       \mathsf N_q,\label{eq:k11}\\
 K_{12}&=-\sum_{e,q}w_q\chi_q\mathsf N_q^T\bn_q\mathsf M_q,
 \label{eq:k12}\\
 K_{21}&=-\sum_{e,q}w_q\chi_q\mathsf M_q^T\bn_q^T\mathsf N_q,
 \label{eq:k21}\\
 K_{22}&=\gamma\sum_{e,q}w_q(\chi_q-1)\mathsf M_q^T\mathsf M_q-S.
 \label{eq:k22}
\end{align}
Before global insertion the contact block dimensions are
$dn_b\times dn_b$, $dn_b\times n_m$, $n_m\times dn_b$, and
$n_m\times n_m$. The jump matrix is assembled separately between
neighboring multiplier DOFs. For example, the scalar-index entry of the
displacement contact block is
\begin{equation}
 (K_{uu}^c)_{(a,i),(b,j)}
 =\sum_q w_q N_a(X_q)N_b(X_q)
       \left(\frac{\chi_q}{\gamma}n_{q,i}n_{q,j}
                        -p_qH_{q,ij}\right).
 \label{eq:indexentry}
\end{equation}
The matching mixed entry is
$-(\sum_qw_q\chi_qN_a n_{q,i}M_e)$, with its transpose in the
other row before scaling. This accounts for both test/trial shape factors
even though the local P0 factors are both the number one.

If $K_{\mathrm{bulk}}$ is symmetric, these unscaled Galerkin blocks give
a symmetric generalized Jacobian: $K_{21}=K_{12}^T$, $\bH$ is symmetric,
and $S=S^T$. Symmetry does not imply positive definiteness. The multiplier
block has a negative inactive mass term, the active system has saddle
coupling, and the curvature term can be negative in tangential directions.
For the sphere, $\bH=(\bI-\bn\bn^T)/r$ makes that last fact explicit.
Zero friction does not imply a zero tangential derivative of a force whose
normal rotates.

For a plane, $\bH=0$. On a strictly inactive point, $p=\chi=0$, so
all displacement/contact couplings vanish and the local multiplier block
is negative $\gamma$ times its mass. On a fully active point, the
local multiplier mass derivative vanishes, but the couplings and
augmentation stiffness remain. At $\lambda=0$ and negative gap,
$p=-g/\gamma$, and the displacement contact tangent reduces to the
separate penalty expression
$\gamma^{-1}(\bn\bn^T+g\bH)$. For nonzero multiplier, replacing
$-p\bH$ by $(g/\gamma)\bH$ would drop a real curvature contribution.

For finite-strain bulk mechanics, $K_{\mathrm{bulk}}$ in
\eqref{eq:k11} must include the full derivative of $\bP$ with respect
to placement. If implemented in second Piola--Kirchhoff form, that
includes the geometric term arising from $\bP=\bF\bS$ as well as
the constitutive derivative. Contact does not supply or repair missing
bulk tangent terms.

\section{What P0 equations enforce, and a hand-integrated edge}
\label{sec:moments}

\subsection{Face moments, not pointwise constraints}
For a contact face $e$, let $A_e$ and $I_e$ denote its active and inactive
quadrature points. Equation~\eqref{eq:vectorrl} at a root reads
\begin{equation}
 -\sum_{q\in A_e}w_qg_q
 -\gamma\Lambda_e\sum_{q\in I_e}w_q
 -(S\Lambda)_e=0.
 \label{eq:facemoment}
\end{equation}
Thus a fully active unstabilized face enforces zero weighted mean gap,
not zero gap at every point. With stabilization its integrated gap equals
$-(S\Lambda)_e$. A fully inactive unstabilized face has $\Lambda_e=0$
when its mass $m_e=\sum_{q\in e}w_q$ is positive; with jumps it instead
satisfies $\gamma m_e\Lambda_e+(S\Lambda)_e=0$.
A partially active face obeys the combined equation
\eqref{eq:facemoment}, not one active flag attached to its P0 unknown.

These distinctions explain why $p_h$ and $\lambda_h$ must be reported
separately. In particular, an inactive face adjacent to loaded faces may
have a nonzero multiplier while its applied pressure remains zero. This
is not adhesive traction. The converged discrete equation does not impose
the pointwise equality $p_h=\lambda_h$ used in the ideal complementarity
argument. Nor does increasing the displacement degree automatically add
multiplier moments when the boundary mesh and P0 space are unchanged.

On a fully active face with no jump term, expanding the contact potential
gives the useful identity
\begin{equation}
 \mathcal L_{c,e}=-\Lambda_e\sum_q w_qg_q
                 +\frac1{2\gamma}\sum_qw_qg_q^2.
 \label{eq:activepotential}
\end{equation}
For positive quadrature weights, write the weighted face average as
$\overline g_e=m_e^{-1}\sum_qw_qg_q$. Then
\begin{equation}
 \sum_qw_qg_q^2=m_e\overline g_e^{\,2}
                 +\sum_qw_q(g_q-\overline g_e)^2.
 \label{eq:gapvariance}
\end{equation}
Even after the multiplier equation removes the mean, the within-face
gap variation remains penalized by $1/\gamma$. P0 is therefore not an
automatic cure for augmentation-driven overconstraint. This derivation
assumes the fully active branch and positive weights; it is not an observed
locking result or a proof of a particular limiting error.

\subsection{Two-node straight edge}
Consider a horizontal straight contact edge of reference length $L$ in
2D, a plane normal $\bn=\boldsymbol e_y$, and a linear displacement
trace. Set $s\in[0,L]$, $N_1=1-s/L$, $N_2=s/L$ and retain only its
two normal displacement coefficients $U_1,U_2$ and one P0 multiplier
$\Lambda$. Assume the edge remains fully active, so
\begin{equation}
 g(s)=g_0(s)+N_1(s)U_1+N_2(s)U_2,\qquad
 p(s)=\Lambda-g(s)/\gamma>0.
 \label{eq:edgegap}
\end{equation}
The plane decomposition is exact. Without bulk, stabilization, or row
scaling, its residuals are
\begin{equation}
 r_i=-\int_0^L p(s)N_i(s)\,ds
     =-\frac L2\Lambda+\frac1\gamma\int_0^L g(s)N_i(s)\,ds,
 \quad i=1,2,\qquad
 r_\Lambda=-\int_0^L g(s)\,ds.
 \label{eq:edgeresidual}
\end{equation}
Using
\begin{equation}
 \int_0^L N_i\,ds=\frac L2,\qquad
 \int_0^L N_i^2\,ds=\frac L3,\qquad
 \int_0^L N_1N_2\,ds=\frac L6,
 \label{eq:edgeintegrals}
\end{equation}
the exact local generalized Jacobian in ordering $(U_1,U_2,\Lambda)$ is
\begin{equation}
 \boxed{K_e=\begin{bmatrix}
 L/(3\gamma)&L/(6\gamma)&-L/2\\
 L/(6\gamma)&L/(3\gamma)&-L/2\\
 -L/2&-L/2&0
 \end{bmatrix}.}
 \label{eq:edgematrix}
\end{equation}
The off-diagonal displacement entry requires the product of two trace
shapes; a lumped or single-point rule would not reproduce this matrix.
The last row couples the \emph{average} normal displacement to the
multiplier. The first two rows also penalize the nonconstant normal mode,
consistent with \eqref{eq:gapvariance}. Tangential displacement entries
are zero for this straight plane, but bulk elasticity still contributes
to both displacement components in a body solve.

If instead the whole edge is strictly inactive, $p=0$ and its contact
residual is $(0,0,-\gamma L\Lambda)^T$. Its only nonzero local tangent
entry is
\begin{equation}
 (K_e)_{\Lambda\Lambda}=-\gamma L.
 \label{eq:edgeinactive}
\end{equation}
Equations~\eqref{eq:edgematrix} and \eqref{eq:edgeinactive} are analytic
calculations from the stated discrete law, not outputs from a newly run
example or a full contact patch test.

\subsection{Multiplier control is a global issue}
The zero active multiplier diagonal in \eqref{eq:edgematrix} makes
normal displacement coupling important. For two equal active linear
collinear edges contacting the same plane with the same admissible normal,
fix the two outer normal DOFs and leave
only the middle normal DOF free. Each multiplier couples to that one
DOF with magnitude $L/2$. With $\delta=0$, the alternating multiplier
vector $(1,-1)^T$ then lies in the null space of the free-displacement
row $K_{12}=-(L/2)\begin{bmatrix}1&1\end{bmatrix}$.
Both edges individually retain a free normal variation, yet their
combined coupling does not control every multiplier mode.

The positive jump matrix controls this particular alternating mode but
leaves the constant mode untouched. The primal coupling must still control
constants and any other unpenalized modes. Neither the P0 count nor a
per-face free-normal check proves a global inf-sup condition. Likewise,
the compliance dimensions and the factor $\delta\gamma h_F$ do not
transfer the scalar source's stability analysis to curved vector contact
on arbitrary meshes. Parameter robustness, locking avoidance, and error
rates need separate analysis and controlled numerical studies.

\section{Monolithic Newton solution and operator composition}
\label{sec:solver}

\subsection{Simultaneous unknowns and accepted load steps}
With displacement-only primal physics, write $Z=(U,\Lambda)^T$.
At a fixed load or prescribed-displacement stage, a semismooth Newton
iteration solves
\begin{equation}
 \begin{bmatrix}K_{11}&K_{12}\\K_{21}&K_{22}\end{bmatrix}_{Z^k}
 \begin{bmatrix}\Delta U\\\Delta\Lambda\end{bmatrix}
 =-\begin{bmatrix}r_U\\r_\Lambda\end{bmatrix}_{Z^k},\qquad
 Z^{k+1}=Z^k+\alpha_k\Delta Z.
 \label{eq:newton}
\end{equation}
$\alpha_k=1$ is a full Newton step; a solver may reduce it if an
appropriate globalization is enabled. There is no staggered
displacement solve followed by a separate multiplier or Uzawa update.
Both fields are current total unknowns. Every evaluation recomputes
activity from the trial vector; no committed active set is hidden in
the contact operator.

An accepted load step supplies an initial guess for the next stage, not
a replacement for the total placement in the gap. A failed increment
must restore both displacement and multiplier coefficients, the accepted
load parameter, and any stateful primal histories. The contact wrapper's
\code{GetChildOperators} exposes its primal operator to the repository's
solver lifecycle traversal. A custom external driver must still implement
the required begin/commit/rollback behavior; calling a residual is not a
commit operation. All trial distance evaluations must remain in smooth
geometry regions. Invalid-geometry assertions may abort rather than
provide a recoverable load cutback.

\subsection{Generic primal extraction and essential masking}
The wrapped primal may contain fields besides displacement. Let
$Y\in\mathbb R^{n_y}$ be its coefficient vector,
$E\in\mathbb R^{n_u\times n_y}$ the fixed linear extraction $U=EY$,
and $R_0(Y)$ its contact-free residual with tangent $J_0(Y)$.
$E$ is not a test field or Young's modulus. Its transpose inserts contact
forces into the corresponding primal coordinates.

Let $Q_D$ be the diagonal displacement coefficient mask that retains
free components. Denote the \emph{contact-only} displacement residual
from \eqref{eq:vectorru} by $r_u^c$, and the unscaled contact matrices
by $K_{uu}^c,K_{12},K_{21},K_{22}$. For positive row weights
$s_u,s_\lambda$, set
\begin{equation}
 D_s=\operatorname{diag}(s_u I_{n_y},s_\lambda I_{n_\lambda}).
 \label{eq:rowscaling}
\end{equation}
The returned residual and generalized Jacobian are
\begin{align}
 \widehat R(Y,\Lambda)
 &=D_s\begin{bmatrix}
       R_0(Y)+E^TQ_Dr_u^c(EY,\Lambda)\\
       r_\Lambda(EY,\Lambda)
       \end{bmatrix},\label{eq:extractionresidual}\\
 \widehat J
 &=D_s\begin{bmatrix}
       J_0+E^TQ_DK_{uu}^cQ_DE&E^TQ_DK_{12}\\
       K_{21}Q_DE&K_{22}
       \end{bmatrix}.
 \label{eq:extractionjacobian}
\end{align}
Prescribed coefficient \emph{values} are retained in $EY$ for gap
evaluation; only residual test rows and Jacobian variations are masked.
Equation~\eqref{eq:extractionjacobian} is the derivative on the essential
constraint manifold, where $(I-Q_D)E\Delta Y=0$, not an unrestricted
derivative with respect to prescribed data. The primal operator and solver
must consistently enforce the prescribed values and their essential
rows. Contact alone does not insert essential identity equations.

The implementation checks matching essential lists for a direct
\code{NonlinearForm} with identity extraction. More general primal
operators and extractions must honor the same contract themselves.
Every selected face must retain a free displacement trace DOF; during
gradient assembly a fully sampled active face is also checked for remaining
constant-mode normal coupling after masking, evaluated directly in
\code{mElementConstantNormalCoupling}. These restrictive local checks are
neither a global rank/inf-sup certificate nor necessary for every stabilized
system to be nonsingular.

Positive row weights preserve the zero set because $D_s$ is invertible.
They change the residual stopping norm and may improve or worsen
conditioning; no improvement follows from positivity alone. Unequal
weights generally destroy Euclidean symmetry, even when the unscaled
system is symmetric. Scaling includes the complete primal row and the
complete multiplier row, including $-S\Lambda$ and $-S$, exactly once.

\subsection{Residual evaluation versus solution, and the circle example}
\code{Mult} evaluates \eqref{eq:extractionresidual} at the supplied vector.
It does not perform Newton, solve for a multiplier, or update unknowns.
\code{GetGradient} assembles the four contact blocks and refreshes the
wrapped primal gradient for \eqref{eq:extractionjacobian}. The returned
block operator and borrowed primal gradient reuse storage; another gradient
evaluation invalidates a retained gradient reference. The global action
uses sparse contact blocks and operator products, not a dense global
matrix or a global contact energy minimization.

The mixed branch of \code{examples/contact/circle.cpp} uses GMRES with a
block-diagonal preconditioner inside
\code{MultiNewtonAdaptive<NewtonLineSearch>}. The primal preconditioner
is a copied, row-scaled bulk elasticity matrix, solved by UMFPACK when
SuiteSparse is available or treated by a symmetric Gauss--Seidel smoother
otherwise. It is not a factorization of the entire mixed Jacobian, nor
does it include the full changing contact augmentation. The multiplier
preconditioner uses a positive diagonal with entries
\begin{equation}
 d_{\lambda,e}=s_\lambda\gamma h_c
 \label{eq:preconditioner}
\end{equation}
on the uniform 2D contact edges. This is the magnitude of the inactive
P0 mass coefficient, not the exact active Schur complement or the
stabilized multiplier block. GMRES accommodates the scaled indefinite
operator; conjugate gradients is not justified for this saddle system.
The example intentionally retains default P0, this preconditioner, and its
one-multiplier-per-face count check; it is not a generic higher-degree solver.

In this example the dimensionless CLI factor $\alpha$ sets
$\kappa=\alpha E_Y/h_c$, where $E_Y$ is Young's modulus, and the mixed
constructor receives
\begin{equation}
 \gamma=\kappa^{-1}=\frac{h_c}{\alpha E_Y},\qquad \delta=1.
 \label{eq:examplegamma}
\end{equation}
The legacy flag \code{-gamma/--penalty-factor} names $\alpha$, not
the paper compliance. The current mixed API takes $\gamma$ in the
former positional stiffness slot: passing the old numerical stiffness
unchanged would alter the model. The separate displacement-only penalty
integrator still takes $\kappa$.

The example scales the primal and multiplier rows by inverse reference
residuals $E_Y u_{\max}W/H_b$ and $u_{\max}W$, respectively, for body
width $W$ and height $H_b$. It applies bottom vertical displacement at
trial pseudo-time while constraining left horizontal displacement
component-wise. Adaptive load increments operate on the complete mixed
vector. The example does not enable the optional line search; the use
of a class bearing that name is not a claim of a proved contact merit
globalization. The fixed contact rule has default order 127 and independent
gap sampling includes face endpoints. These are example resolution choices,
not a collision-detection bound or a library-wide quadrature guarantee.

\section{Verification map and limits of interpretation}
\label{sec:verification}

\subsection{Existing source checks}
The following are inspected test constructions in
\code{tests/contact_test.cpp}, not newly executed results. Unless noted,
the suite prefix is \code{SemismoothContact}. Long identifiers are given
verbatim so the derivation can be traced to a specific regression.

\testname{ActivePlaneResidualHasExpectedPrimalAndMultiplierResultants}
Checks the active-plane residual sign and magnitude, zero tangential
resultant, positive physical obstacle resultant, and the integrated gap
in the multiplier row. Its companion
\code{InactivePlaneConstrainsMultiplierWithoutDisplacementForce}
checks zero displacement force and the negative compliance-scaled
multiplier mass. \code{SupportsThreeDimensionalPlaneContact} checks
reference area factors and a mixed directional derivative in 3D.

\testname{GammaAndDeltaMatchIntegratedTwoFaceReference}
Hand-integrates a two-face planar reference using a constant normal
displacement direction and two P0 multiplier directions. It compares
residuals and every projected matrix block for active/inactive cases,
multiple compliance and jump-factor values, and unequal row scales.
\code{RejectsInvalidComplianceAndJumpWeight} checks invalid parameters,
including a positive compliance whose reciprocal overflows.

\testname{CurvedMixedJacobianMatchesDirectionalDifferenceWithExtraPrimalField}
Compares the full generalized Jacobian action with a centered residual
difference on a curved active branch, including unequal row weights and
three additional primal DOFs. This exercises the normal derivative and
extraction without identifying extra fields as displacement.
\code{EqualityUsesTheActiveOneSidedLinearization} instead uses an
active-side perturbation at equality, as required by the selected kink
derivative.

\testname{HighOrderBoundaryTraceMatchesVolumeOracle}
Compares trace assembly with an independent adjacent-volume-shape oracle
for curved high-order triangles, quadrilaterals, tetrahedra, and hexahedra,
both global orderings, and a nonsymmetric custom quadrature rule. It
also checks trace-zero directions. Stabilization is disabled in this
oracle; jump behavior is covered separately.

\testname{BoundaryP0JumpStabilizationCouplesAdjacentFaces}
Checks the jump sign, conservation of the sum of jump residuals, and an
alternating multiplier tangent action on inactive faces.
\code{CurvedThreeDimensionalJumpUsesIntegratedInterfaceMeasure}
compares the stabilized-minus-unstabilized response against independently
integrated curved edge measure. These are coefficient and assembly checks,
not active-pair inf-sup tests.

\testname{EssentialDisplacementDofsAreRemovedFromContactVariations}
Checks eliminated contact directions while retaining their prescribed
values in the gap. The guards
\code{RejectsFullyActiveFaceWithoutFreeNormalVariation} and
\code{RejectsEssentialDofsThatDifferFromPrimalForm} cover corresponding
input contracts. \code{BoundaryP0HasOneMultiplierPerSelectedFace} and
\code{WrapsBlockNonlinearPrimalOperator} cover P0 count and a composite
primal with an unrelated scalar field. They do not prove stability of
arbitrary extractions or multiphysics couplings.

\subsection{Algebraic residual, sampled contact error, and missed contact}
The diagnostic contact resultant is
\begin{equation}
 \boldsymbol f_c^Q=\sum_qw_qp_q\bn_q,
 \qquad e_{\mathrm{map}}^Q=\gamma\max_q|p_q-\lambda_q|,
 \qquad e_{\mathrm{pen}}^Q=\max_q\max(0,-g_q).
 \label{eq:diagnostics}
\end{equation}
The latter two quantities have length units. The sampled contact-map
error is not the assembled multiplier residual norm: projection and jumps
allow it to remain nonzero at a discrete root. The minimum sampled gap
is not the minimum over the continuous deformed trace.

\code{CoarseQuadraticCircleNeedsResolvedQuadrature} constructs a
two-by-two quadratic square mesh whose six default assembly samples
miss a penetrating circle patch. Its prescribed upward translation is
$0.02$, circle center $(0.65,1.26)$, and radius $0.25$; independent gap
sampling detects a minimum of $-0.01$. These are the test's prescribed
construction and analytic geometry, not a new simulation report. The
test then checks resolved resultant integration and branchwise
directional derivatives with higher-order rules.
\code{IndependentGapSamplingUsesCurrentCurvedGeometry} separately
checks evaluation of the deformed curved trace, including endpoints.

Accordingly, successful Newton termination should be accompanied by
separate reporting of reaction/resultant, sampled penetration,
unprojected contact-map error, and quadrature/sampling sensitivity.
Independent finite sampling improves detection but still cannot certify
nonpenetration between samples. The circle example's permissive
penetration/contact-map acceptance threshold is tied to the imposed
displacement, so its acceptance is a smoke/resolution check rather than
a tight pressure-accuracy criterion.

\subsection{Supported scope and unresolved validation}
The current contact implementation is serial, conforming, default-CPU,
and based on fixed reference measures with analytic smooth rigid walls.
It rejects parallel displacement spaces and nonconforming meshes and
requires supported H1 traces without DOF transformations. It has no MPI
contact assembly, GPU kernel, two-body search, self-contact, friction,
or contact-history transfer. Borrowed primal, extraction, displacement
space, obstacle, integration rule, and boundary multiplier object must outlive
the operator. The boundary object owns its submesh geometry, collection, and
FE space, not the solution, and borrows the parent mesh. Both it and contact
prohibit copy/move. Scratch is mutable and not reentrant. Cached maps and jump
matrices require reconstruction of the boundary object and its borrowers after
mesh, geometry, or space changes. NURBS and curved \code{SubMesh} parents are
unsupported; curved parent nodes must be uniform-order Gauss--Lobatto H1.
Discontinuous geometry is rejected because MFEM's coefficient-copy geometry
transfer can duplicate parent DOFs. Multiplier extrema
are quadrature-sampled field values, not coefficient extrema or a between-sample
positivity certificate.
The wrapped primal must be contact-free to avoid counting the same
interaction twice.

Neither the source contact algebra nor the tests listed above establishes
existence, uniqueness, mesh-independent inf-sup stability, locking
avoidance, or a convergence rate for this vector formulation. No new
flat-contact patch pressure-error study, curved reference-solution study,
mesh-to-mesh comparison, or controlled parameter sweep is supplied here.
Such studies should separate bulk approximation, geometry, multiplier
resolution, quadrature, and nonlinear tolerances. Decreasing $\gamma$
at fixed mesh strengthens augmentation while weakening the stated jump
coefficient; it is not interchangeable with a justified mesh-dependent
refinement sequence. An exact residual derivative is necessary for a
reliable Newton method, but is not by itself a discretization theorem.

\section{Short code map and closing perspective}
\label{sec:codemap}

The central implementation is \code{plugin::SemismoothRigidContactOperator}
in \code{fem/Contact.h} and \code{fem/Contact.cpp}. Private methods and state
are declared directly in the operator's header, with method definitions in
the source file. Source locators use names rather than line numbers so that
routine edits do not invalidate the map.

\begin{table}[!ht]
\centering
\small
\caption{Mathematical quantities and their implementation locations.}
\label{tab:codemap}
\begin{tabular}{@{}p{0.45\linewidth}p{0.50\linewidth}@{}}
\toprule
Location & Mathematical role\\
\midrule
\code{RigidPlaneObstacle::Evaluate}; \code{RigidSphereObstacle::Evaluate}
 & $g,\bn,\bH$ in \eqref{eq:distance}, \eqref{eq:planegap}, and
 \eqref{eq:sphere}.\\[4pt]
\code{VerifyInput}; \code{BoundaryMultiplierSpace::VerifyElementMaps}
 & Parameters, fields, constraints; boundary ownership and verified local
 vertex/DOF maps in \code{fem/BoundaryMultiplierSpace.cpp}. Contact uses
 \code{GetParentBoundaryElement} and \code{GetElementDofs}.\\[4pt]
\code{VisitElementQuadraturePoints}; \code{BuildIntegrationRules}
 & Trace interpolation, orientation, $\bx_q$, $w_q$, $p_q$, and
  $\chi_q$ in \eqref{eq:quadrature}--\eqref{eq:pointkernel}. Per-element rule
  pointers are cached at construction and rebuilt by \code{SetIntegrationRule};
  residual, Jacobian, and diagnostics reuse them. Reconstruct after
  mesh/geometry/space changes.\\[4pt]
\code{AssembleContact}
 & Explicit element loop: reset, integrate, scatter; then assemble jumps.\\[4pt]
\code{AccumulateContactPointResidual}; \code{AccumulateContactPointJacobian}
 & Contact residuals and the four blocks
 \eqref{eq:k11}--\eqref{eq:k22}, excluding the separately assembled jumps.\\[4pt]
\code{BuildMultiplierJumpStabilization}
 & Direct interface quadrature and cached basis-product jump matrices;
 \eqref{eq:hface}--\eqref{eq:jumpmatrix} is the default P0 reduction.\\[4pt]
\code{AssembleMultiplierJumpStabilization}
 & $-S\Lambda$, $-S$, and the multiplier row scale applied once.\\[4pt]
\code{ScatterContactElement}; \code{MaskEssentialElementEntries}
 & $Q_D$ in \eqref{eq:extractionresidual}--\eqref{eq:extractionjacobian};
 free-normal validation and global insertion. Prescribed values are not
 removed from gap evaluation.\\[4pt]
\code{Mult}; \code{GetGradient}
 & Residual evaluation and generalized Jacobian refresh, respectively;
 neither performs a nonlinear solve.\\[4pt]
\code{ComputeContactDiagnostics}
 & Unscaled quadrature resultant, gap, multiplier, pressure, and
 contact-map diagnostics in \eqref{eq:diagnostics}.\\[4pt]
\code{examples/contact/circle.cpp}
 & Small-strain bulk with nonlinear rigid-circle distance, row scales,
 compliance conversion, GMRES/block preconditioner, adaptive loading.\\[4pt]
\code{fem/Solvers.h}; \code{fem/Solvers.cpp}
 & Nonlinear/load-step drivers and composite primal lifecycle support.\\[4pt]
\code{tests/contact_test.cpp}
 & Named geometry, residual, tangent, stabilization, constraint,
 extraction, and quadrature checks in Section~\ref{sec:verification}.\\
\bottomrule
\end{tabular}
\end{table}
\FloatBarrier

The essential translation is compact: the scalar source contact trace
becomes the negative \emph{total gap}, its variation becomes
$-\bn\cdot\bv$, and compression reverses the source multiplier sign.
The rest follows from consistent differentiation and discretization.
Initial clearance must not be discarded, the curved normal must be
differentiated, the multiplier row must retain its saddle normalization,
and the P0 equation must be interpreted as a stabilized moment.
Those distinctions connect the prominent vector pair
\eqref{eq:counterpart} to the actual discrete system
\eqref{eq:finalweak} without attributing more to either the source paper
or the implementation than has been established.

\begingroup
\small
\begin{thebibliography}{9}
\bibitem{BHLv1}
E. Burman, P. Hansbo, M. G. Larson,
Augmented Lagrangian finite element methods for contact problems,
arXiv:1609.03326v1, 12 September 2016, 26 pp.
\url{https://arxiv.org/abs/1609.03326v1};
versioned PDF: \url{https://arxiv.org/pdf/1609.03326v1}.

\bibitem{BHLjournal}
E. Burman, P. Hansbo, M. G. Larson,
Augmented Lagrangian finite element methods for contact problems,
ESAIM: Mathematical Modelling and Numerical Analysis 53 (1) (2019)
173--195.
\url{https://doi.org/10.1051/m2an/2018047}.
\end{thebibliography}
\endgroup
\end{document}
